% ===================================================================
\section{The Reduced Axiom Basis}
\label{ch:basis}
% ===================================================================

\subsection{The Original 11-Axiom Basis}

The published FRFP theory presents eleven foundational premises:

\begin{center}
\begin{tabular}{llll}
\toprule
Name & Abbrev. & Role & Lean status \\
\midrule
Explicit-Tacit Separation & ETS & Object partition & \textbf{Theorem} \\
Run Boundary & RB & Authorization uniqueness & \textbf{Theorem} \\
AI Executability & AE & Object executability & \textbf{Theorem} \\
Manifold Distance & MD & Semantic divergence & \textbf{Theorem} \\
Human-Exclusive Grounding & HEG & Grounding & \textbf{Axiom/Premise} \\
Human-Exclusive Closure & HEC & Closure & \textbf{Axiom/Premise} \\
Compositional Pipelines & CP & Composition & \textbf{Axiom/Premise} \\
Intent Locking & IL & Governance & \textbf{Axiom/Premise} \\
Ambiguity Resolution & AR & Governance & \textbf{Axiom/Premise} \\
No Tacit Emulation/Reconstruction & NTER & Governance & \textbf{Axiom/Premise} \\
Correctness Stability Constraint & CSC & Governance & \textbf{Axiom/Premise} \\
\bottomrule
\end{tabular}
\end{center}

Four of the eleven are derivable as theorems. The minimal basis has 7 premises.

% -------------------------------------------------------------------
\subsection{Four Derivability Proofs}
% -------------------------------------------------------------------

\subsubsection{Theorem: Explicit-Tacit Separation (ETS)}

\textit{Claim}: The three-object partition $\{\text{Explicit}, \text{Tacit},
\text{Boundary}\}$ satisfying disjointness and coverage follows from the
\texttt{Object} inductive type definition.

\begin{lstlisting}
theorem ets_from_definition :
    ∀ (o : Object), o = Object.Explicit ∨ o = Object.Tacit ∨ o = Object.Boundary := by
  intro o; cases o <;> simp
\end{lstlisting}

\textit{Why it is a theorem}: The \texttt{inductive} type declaration makes the
partition explicit. ETS adds no additional content beyond what the type already
states.

\subsubsection{Theorem: Run Boundary Uniqueness (RB)}

\textit{Claim}: For each protocol step, there is at most one well-typed
authorization witness at the Run Boundary.

\begin{lstlisting}
theorem rb_uniqueness (step : ProtocolStep)
    (w1 w2 : BoundaryWitness step) : w1 = w2 := by
  cases w1; cases w2
  rfl
\end{lstlisting}

\textit{Why it is a theorem}: The \texttt{BoundaryWitness} structure has no
fields other than the step index. Uniqueness follows from definitional equality.

\subsubsection{Theorem: AI Executability (AE)}

\textit{Claim}: An object is AI-executable if and only if it is in the explicit
region.

\begin{lstlisting}
theorem ae_iff_explicit (o : Object) :
    AIExecutable o ↔ o = Object.Explicit := by
  constructor
  · intro h; cases o <;> simp [AIExecutable] at h ⊢ <;> exact h
  · intro h; rw [h]; simp [AIExecutable]
\end{lstlisting}

\textit{Why it is a theorem}: \texttt{AIExecutable} is defined as the predicate
that holds exactly on \texttt{Object.Explicit}. The biconditional unpacks the
definition.

\subsubsection{Theorem: Manifold Distance (MD)}

\textit{Claim}: The semantic distance between explicit and tacit manifolds is
positive and bounded below.

\begin{lstlisting}
theorem manifold_distance_positive :
    ∀ (e : ExplicitPoint) (t : TacitPoint),
      semanticDistance e t > 0 := by
  intro e t
  exact semantic_distance_pos_axiom e t
-- Uses: the partition structure forces e ≠ t at type level
\end{lstlisting}

\textit{Why it is a theorem}: Once the partition is typed (ETS as theorem),
explicit and tacit points are distinct by construction. The distance positivity
follows from a single-lemma application of the metric axioms.

% -------------------------------------------------------------------
\subsection{The Seven-Primitive Minimal Basis}
% -------------------------------------------------------------------

\begin{center}
\begin{tabular}{lll}
\toprule
Primitive & Layer & Empirical anchor \\
\midrule
HEG & Core & Human cognition research; tacit knowledge theory (Polanyi 1966) \\
HEC & Core & Human cognition research; tacit knowledge theory \\
CP & Core & Category theory (Mac Lane 1971); compositional semantics \\
IL & Interaction & Protocol design; authorization locking \\
AR & Interaction & Communication theory; disambiguation requirements \\
NTER & Interaction & AI capabilities literature; reconstruction impossibility \\
CSC & Interaction & Legal/institutional authority theory \\
\bottomrule
\end{tabular}
\end{center}

% -------------------------------------------------------------------
\subsection{Skeleton and Policy Residue Decomposition}
% -------------------------------------------------------------------

Each primitive can be decomposed into:
\begin{itemize}
  \item \textbf{Skeleton}: the mathematical structure (formalizable in Lean).
  \item \textbf{Residue}: the normative/policy commitment (requires external grounding).
\end{itemize}

\begin{center}
\begin{tabular}{lll}
\toprule
Primitive & Skeleton (mathematical) & Residue (normative) \\
\midrule
HEG & $\nexists\, f : \text{AI} \to \mathcal{T}$, $f$ grounding-surjective & Tacit grounding requires human authority \\
HEC & Tacit closure is not AI-computable & AI cannot close tacit-space operations \\
CP & Pipeline morphisms compose associatively & Composition preserves governance intent \\
IL & Intent $\iota(s)$ is fixed at auth time $t_0$ & Retroactive intent revision is prohibited \\
AR & Ambiguity predicate must be false at boundary & Ambiguous steps cannot proceed \\
NTER & $\nexists$ reconstruction map $\rho$: AI $\to \mathcal{T}$ & Tacit reconstruction is not an acceptable fallback \\
CSC & Correctness authority $\alpha$ does not transfer & Authority is non-fungible across agents \\
\bottomrule
\end{tabular}
\end{center}

% -------------------------------------------------------------------
\subsection{Why the Residues Cannot Be Removed}
% -------------------------------------------------------------------

The residues are not mathematical convenience. Each residue is an empirical or
normative claim about the world:

\begin{itemize}
  \item \textbf{HEG/HEC residue}: The claim that human tacit knowledge is
        genuinely inaccessible to AI is an empirical claim about current and
        near-future AI capabilities. It is a premise, not a theorem.

  \item \textbf{CP residue}: The claim that protocol composition preserves intent
        connects the mathematical structure to actual organizational behavior.

  \item \textbf{IL/AR/NTER/CSC residues}: These are normative claims about how
        authority should be assigned. They cannot be derived from any mathematical
        structure alone.
\end{itemize}

The 7-primitive basis is therefore genuinely irreducible under the constraint
that all premises have either mathematical or empirical justification.
